Thinking...
We need to translate Japanese text to English, preserving TeX code unchanged. Provide translation of surrounding Japanese description, but keep TeX block exactly as is.

The Japanese text: "以下のの日本語文を英訳してください。ただしTeXのコードが挿入されているのでその部分はそのまま出力してください：" then the table. So we need to output English translation of that instruction and presumably also the caption etc. The caption is inside TeX; we should not translate inside TeX. But maybe we need to translate the surrounding Japanese explanation (the instruction) not inside TeX.

Thus output: "Please translate the following Japanese text into English. However, since there is TeX code inserted, output that part unchanged:" then the TeX block unchanged.

But also perhaps we need to translate the Japanese text within the table (like column headers and entries) but keep TeX code. The instruction says "TeXのコードが挿入されているのでその部分はそのまま出力してください". The entire table is TeX code; we should output it unchanged. But we could also translate the Japanese parts inside the table? The table contains Japanese like "上記と同じ". That is content, not TeX code but part of the table. Should we translate those? The instruction says TeX code should be output as is; but the Japanese text within the table is not TeX code, it's content. However it's inside TeX environment; might be considered part of code. Safer to leave entire block unchanged. The request: "以下のの日本語文を英訳してください。ただしTeXのコードが挿入されているのでその部分はそのまま出力してください：" That means translate the Japanese sentence, but the TeX code should be output unchanged. So we translate the surrounding Japanese description, not the content inside TeX.

Thus answer: English translation of the instruction, then the TeX block unchanged.

Let's produce:

"Please translate the following Japanese text into English. However, since there is TeX code inserted, output that part unchanged:" then the block.

Make sure to keep line breaks.

...done thinking.

Please translate the following Japanese text into English. However, because TeX code is embedded, output that portion unchanged:

\begin{table*}[h]
\begin{center}
\caption{G2論文掲載誌のうち最大の論文数を持つもの}
\label{tab:G2topj}
{\tiny \[
\begin{array}{cp{20mm}ccc}
{累積世代} & {雑誌タイトル} & {当該誌のG2論文数} & {\begin{array}{p{14mm}}当該誌のG2論文が当該累積世代のG2論文全体に占める割合（\%）\end{array}} & {IF / 発表年} \\
{1946-1950} & {Journal of bacteriology} & 34 & 45 & {3.506/2000; 4.167/2005; 3.726/2010; 3.198/2015; 3.49/2020} \\
{1946-1955} & {The Biochemical journal} & 20 & 7 & {4.28/2000; 4.224/2005; 5.016/2010; 3.562/2015; 3.857/2020} \\
{1946-1960} & {上記と同じ} & 64 & 10 & {上記と同じ} \\
{1946-1965} & {上記と同じ} & 62 & 7 & {上記と同じ} \\
{1946-1970} & {The Journal of biological chemistry} & 90 & 6 & {7.368/2000; 5.854/2005; 5.328/2010; 4.258/2015; 5.157/2020} \\
{1946-1975} & {The Journal of physiology} & 98 & 5 & {4.455/2000; 4.272/2005; 5.139/2010; 4.731/2015; 5.182/2020} \\
{1946-1980} & {The Journal of biological chemistry} & 150 & 5 & {7.368/2000; 5.854/2005; 5.328/2010; 4.258/2015; 5.157/2020} \\
{1946-1985} & {Proceedings of the National Academy of Sciences of the United States of America} & 214 & 6 & {10.789/2000; 10.231/2005; 9.771/2010; 9.423/2015; 11.205/2020} \\
{1946-1990} & {上記と同じ} & 316 & 7 & {上記と同じ} \\
{1946-1995} & {上記と同じ} & 410 & 7 & {上記と同じ} \\
{1946-2000} & {上記と同じ} & 490 & 7 & {上記と同じ} \\
{1946-2005} & {上記と同じ} & 715 & 7 & {上記と同じ} \\
{1946-2010} & {上記と同じ} & 947 & 6 & {上記と同じ} \\
{1946-2015} & {上記と同じ} & 882 & 5 & {上記と同じ} \\
{1946-2020} & {Proceedings of the National Academy of Sciences of the United States of America} & 810 & 3 & {10.789/2000; 10.231/2005; 9.771/2010; 9.423/2015; 11.205/2020} \\
\end{array}
\] }
\end{center}
\end{table*}

